% Bewertungspaket zum Gesamttest «Dynamik» — Quelle des PDFs (python3 scripts/build-lp-pdf.py dynamik).
% Ganze Punkte je Teilschritt; Werte und Folgefehler-Fälle mit python3 nachgerechnet (04.10.2026).
% Fassung 3 (06.10.2026): G1d, G2–G6 neu (keine Wiederholung von Kapitel-, Leisten- oder Übungsaufgaben),
% gleiche Fehler gleich bewertet, Raster zeilenweise mit mehr Luft; Werte und Folgewerte mit python3 nachgerechnet.
% Aufbau, Auftrag an die KI und Punktarten (E, W, B, A) wie im Leitprogramm Kinematik.
\documentclass[11pt]{article}
\usepackage{../lp-druck}
\renewcommand{\lpname}{Dynamik}
\renewcommand{\lpteilgebiet}{4.2}
\renewcommand{\lpfuss}{Bewertungspaket}
\newcolumntype{P}{>{\centering\arraybackslash}p{10mm}}
% arraystretch 1.5: Brüche benachbarter Zeilen berühren sich sonst
\newenvironment{raster}{\par\smallskip\noindent\renewcommand{\arraystretch}{1.5}\tabularx{\linewidth}{@{}>{\raggedright\arraybackslash}p{38mm}P >{\raggedright\arraybackslash}X@{}}\toprule
  \small\textsc{Teilschritt} & \small\textsc{P} & \small\textsc{Musterlösung}\\\midrule}{\bottomrule\endtabularx\par}
\newcommand{\fehler}[1]{\par\smallskip{\small\textcolor{lprot}{\bfseries Typische Fehler:} #1}}
\newcommand{\E}{\,{\footnotesize\bfseries\color{lpgruen}(E)}}
\newcommand{\B}{\,{\footnotesize\bfseries(B)}}
\newcommand{\W}{\,{\footnotesize\bfseries(W)}}
\newcommand{\A}{\,{\footnotesize\bfseries\color{lpgruen}(A)}}
\begin{document}
\lpkopf{Bewertungspaket zum Gesamttest}{}

\begin{lpachtung}{Erst lösen, dann öffnen}
Dieses Paket enthält die Musterlösung. Wer es vor dem Lösen liest, misst am Ende nur, wie gut das Abschreiben gelingt.
\end{lpachtung}

\section*{1 · So gehst du vor}
\begin{enumerate}[leftmargin=*,itemsep=2pt]
\item \textbf{Gesamttest lösen} — vollständig, mit Rechenweg, ohne dieses Paket.
\item \textbf{Fotografieren oder scannen}, gut lesbar, eine Seite pro Bild. \textbf{Kein Name} auf den Bildern.
  Handyfotos tragen oft unsichtbar Standort, Zeit und Gerät mit (Metadaten). Lade darum lieber einen Scan oder ein
  Bildschirmfoto des Fotos hoch, oder entferne vorher die Standortangaben.
\item \textbf{Bewerten lassen.} Ob du dafür eine KI nutzt und welche, entscheidest du selbst und in eigener Verantwortung —
  je nachdem, was dir aufgrund deines Alters und deiner persönlichen Situation erlaubt ist (Altersgrenzen und
  Nutzungsbedingungen des Anbieters, allenfalls Einverständnis deiner Eltern). Lade nur deine Lösung und dieses PDF hoch,
  keine persönlichen Daten, und füge den Auftrag aus Abschnitt~2 ein. Ohne KI kannst du dich mit Abschnitt~3 selbst bewerten.
\item \textbf{Rückmeldung prüfen.} Stimmt, was die KI aus deiner Handschrift gelesen hat? Eine KI kann sich irren —
  im Zweifel gilt das Raster in Abschnitt~3.
\item \textbf{Gezielt wiederholen} nach Abschnitt~4.
\end{enumerate}

\section*{2 · Auftrag an die KI}
\begin{tcolorbox}[breakable,colback=lplinie!25,colframe=lplinie,boxrule=0.5pt,arc=1mm,fontupper=\small]
Du bist Korrektorin oder Korrektor für Physik an einer Schweizer Berufsmaturitätsschule. Im Anhang findest du (1)~das Bewertungspaket zum Gesamttest «Dynamik» mit Musterlösung und Punkteraster und (2)~Fotos meiner handschriftlichen Lösung.\par\medskip
Geh so vor:\par
1. Schreib zu jeder Aufgabe G1 bis G6 zuerst ab, was du in meiner Lösung liest — Rechenweg und Ergebnis. Was du nicht sicher lesen kannst, markierst du mit [unsicher]. Nicht raten.\par
2. Vergib die Punkte genau nach dem Raster im Bewertungspaket (Abschnitt 3), ganze Punkte je Zeile. Ziehe einen Fehler nur einmal ab: Rechne ich mit einem falschen Zwischenergebnis richtig weiter, gibt es die Folgepunkte — auch wenn das Folgeergebnis unplausibel ist; eine fehlende Plausibilitätsprobe kostet keinen zusätzlichen Punkt. Derselbe Fehler in mehreren Teilaufgaben (z.\,B. Taschenrechner im Bogenmass, sin/cos vertauscht, dieselbe fehlende Umrechnung) kostet nur beim ersten Vorkommen einen Punkt; die anderen Teilaufgaben werden folgerichtig bewertet. Gleichartige Fehler kosten gleich viel. Die Zeilen «Typische Fehler» sagen, welche Rasterzeile bei diesem Fehler 0 Punkte gibt — sie sind kein zusätzlicher Abzug.\par
3. Andere richtige Lösungswege zählen voll. Gleichwertige Schreibweisen zählen voll: Dezimalpunkt oder Dezimalkomma, andere Einheit oder Vorsilbe ($0.12\;\text{m} = 12\;\text{cm}$; $\text{m/s}$ statt $\text{km/h}$, wenn richtig umgerechnet), Zehnerpotenz oder ausgeschriebene Zahl, sinnvoll gerundete Werte (drei signifikante Stellen). Zeilen mit (E) sind Ergebnispunkte: Sie verlangen das Ergebnis \emph{mit richtiger Einheit}, auch ohne Weg; ohne oder mit falscher Einheit gibt es diesen Punkt nicht. Zeilen mit (W) gibt es nur mit nachvollziehbarem Weg (Formel, dann Werte mit Einheiten); ein Zahlenwert darin braucht ebenfalls die Einheit. Zeilen mit (B) gehören zu Aufgaben, die eine Begründung, Erklärung oder Formulierung in eigenen Worten verlangen: Es zählt die fachliche Aussage, eine Formel ist nicht nötig; andere richtige Formulierungen zählen voll. Zeilen mit (A) sind Ablesepunkte: der richtig aus dem Diagramm abgelesene Wert mit Einheit. Werte aus sinnvoll gerundeten Zwischenergebnissen zählen, wenn sie höchstens rund 2\,\% vom Raster abweichen.\par
4. Begründe jeden Abzug in einem Satz und nenne die Fehlerart (zum Beispiel «Wurzel vergessen»). Schreib die Musterlösung nicht ab — zeig mir, wo mein Fehler liegt.\par
5. Gib zum Schluss eine Tabelle «Aufgabe | Punkte | Begründung», die Gesamtpunktzahl (von 25) und — nach der Selbsteinschätzung in Abschnitt 4 — welches Kapitel des Leitprogramms ich wiederholen soll. Keine Note.\par
6. Fehlt ein Foto oder ist eine Aufgabe unlesbar, sag es, statt sie zu bewerten.
\end{tcolorbox}

\needspace{8\baselineskip}
\section*{3 · Musterlösung und Punkteraster}
\textbf{Allgemein:} Ganze Punkte je Zeile. Ein Fehler wird nur einmal abgezogen: Wer mit einem falschen Zwischenergebnis
richtig weiterrechnet, bekommt die folgenden Zeilen (Folgefehler) — auch wenn das Folgeergebnis unplausibel ist; eine fehlende
Plausibilitätsprobe kostet keinen zusätzlichen Punkt. Derselbe Fehler in mehreren Teilaufgaben (z.\,B. Taschenrechner im Bogenmass, sin/cos vertauscht, dieselbe fehlende Umrechnung) kostet nur beim ersten Vorkommen einen Punkt; die anderen Teilaufgaben werden folgerichtig bewertet. Gleichartige Fehler kosten gleich viel: ein Fehler, ein Abzug. Gleichwertige Wege und Schreibweisen zählen voll.
In der Spalte P: \textbf{(E)}~= Ergebnispunkt — es zählt das Ergebnis \emph{mit richtiger Einheit}, auch ohne Weg; ohne oder mit
falscher Einheit gibt es ihn nicht. \textbf{(W)}~= Wegpunkt — nur mit nachvollziehbarem Weg (Formel, dann Werte mit Einheiten);
steht in einer Weg-Zeile ein Zahlenwert, braucht auch er die Einheit. \textbf{(B)}~= Begründungspunkt, nur bei Aufgaben, die eine Begründung, Erklärung oder Formulierung verlangen — es zählt die fachliche
Aussage in eigenen Worten, eine Formel ist nicht nötig. \textbf{(A)}~= Ablesepunkt — der richtig abgelesene Wert mit Einheit.
Folgewerte aus gerundeten Zwischenergebnissen zählen, wenn sie höchstens rund 2\,\% abweichen. «Typische Fehler» nennt die Zeilen mit 0 Punkten und die Punktzahl, die bleibt.

\begin{lpaufgabe}{G1}{Begriffe}{4 P}
\begin{raster}
a) Kraft & 1\B & an ihrer Wirkung: Der Körper wird verformt, oder seine Bewegung ändert sich (schneller, langsamer, andere Richtung)\\
b) Trägheitsgesetz & 1\B & Wirkt keine (Gesamt-)Kraft, bleibt ein Körper in Ruhe oder bewegt sich gleichförmig geradeaus weiter (er behält seinen Bewegungszustand)\\
c) Faktor & 1\B & $6$: $a = \dfrac{F}{m}$, Zähler mal $3$, Nenner durch $2$, also $a$ mal $3 \cdot 2 = 6$ (Faktor mit Begründung nötig)\\
d) Kisten & 1\B & Beide gleich stark anstossen (gleiche Kraft): Die Kiste mit der kleineren Beschleunigung hat die grössere Masse, denn $a = \dfrac{F}{m}$ — die Masse ist das Mass für die Trägheit. Gleichwertig: «die volle lässt sich schwerer in Bewegung setzen oder abbremsen».\\
\end{raster}
\fehler{c) Faktor $1.5$ ($3 : 2$) oder $\tfrac16$: c) 0 $\to$ 3 von 4\,P. d) «Man wiegt sie» oder «die schwerere fällt schneller» (im schwerelosen Raum ohne Wirkung): d) 0 $\to$ 3 von 4\,P.}
\end{lpaufgabe}

\begin{lpaufgabe}{G2}{Schlitten: $50\;\text{kg}$, erst $60\;\text{N}$ bei konstant $1.5\;\text{m/s}$, dann $100\;\text{N}$}{3 P}
\begin{raster}
Reibung & 1\B & $F_W = 60\;\text{N}$: konstantes Tempo heisst $a = 0$, also $F_\text{ges} = 0$ und Zug $=$ Reibung (Trägheitsgesetz; Begründung nötig)\\
Beschleunigung & 1\E & $a = \dfrac{F_\text{ges}}{m} = \dfrac{100\;\text{N} - 60\;\text{N}}{50\;\text{kg}} = 0.8\;\text{m/s}^2$\\
Tempo & 1\E & $v = v_0 + a \cdot t = 1.5\;\text{m/s} + 0.8\;\text{m/s}^2 \cdot 3\;\text{s} = 3.9\;\text{m/s}$\\
\end{raster}
\fehler{Reibung $0$ («er gleitet ja von selbst»): Reibung 0; folgerichtig $a = 2\;\text{m/s}^2$ und $v = 7.5\;\text{m/s}$ zählen $\to$ 2 von 3\,P.
$a = \dfrac{100\;\text{N}}{50\;\text{kg}} = 2\;\text{m/s}^2$ trotz richtiger Reibung: Beschleunigung 0; $v = 7.5\;\text{m/s}$ folgerichtig $\to$ 2 von 3\,P.
$v = a \cdot t = 2.4\;\text{m/s}$ (Anfangstempo vergessen): Tempo 0 $\to$ 2 von 3\,P.}
\end{lpaufgabe}

\begin{lpaufgabe}{G3}{Widerstand aus dem Diagramm: $1200\;\text{kg}$, Antrieb $3000\;\text{N}$}{5 P}
\begin{raster}
a) Ablesen & 1\A & $\Delta v = 12\;\text{m/s}$ in $\Delta t = 6\;\text{s}$, also $a = 2\;\text{m/s}^2$ (zählt auch, wenn die Werte direkt in $F = m \cdot \dfrac{\Delta v}{\Delta t}$ stehen)\\
a) Gesamtkraft & 1\E & $F_\text{ges} = m \cdot a = 1200\;\text{kg} \cdot 2\;\text{m/s}^2 = 2400\;\text{N}$\\
b) Widerstand & 1\W & $F_\text{ges} = F_A - F_W$, also $F_W = F_A - F_\text{ges} = 3000\;\text{N} - 2400\;\text{N} = 600\;\text{N}$\\
c) Gesamtkraft beim Ausrollen & 1\E & $a = \dfrac{8\;\text{m/s} - 12\;\text{m/s}}{8\;\text{s}} = -0.5\;\text{m/s}^2$, $F_\text{ges} = 1200\;\text{kg} \cdot (-0.5\;\text{m/s}^2) = -600\;\text{N}$ (Vorzeichen oder «gegen die Fahrtrichtung» nötig)\\
c) gleich gross? & 1\B & Ja: Ohne Antrieb ist der Widerstand die einzige Kraft längs der Fahrt, $F_\text{ges} = -F_W$, also $F_W = 600\;\text{N}$ wie beim Anfahren\\
\end{raster}
\fehler{a) $F = m \cdot v = 14\,400\;\text{N}$ (Geschwindigkeit statt Beschleunigung): a) Gesamtkraft 0; a) Ablesen zählt, wenn $\Delta v = 12\;\text{m/s}$ und $\Delta t = 6\;\text{s}$ richtig dastehen; b) folgerichtig $F_W = 3000\;\text{N} - 14\,400\;\text{N} = -11\,400\;\text{N}$ zählt $\to$ 4 von 5\,P.
b) $F_W = F_A + F_\text{ges} = 5400\;\text{N}$ (Vorzeichen): b) 0 $\to$ 4 von 5\,P; c) «nein, 5400 N gegen 600 N» folgerichtig zählt.
c) $+600\;\text{N}$ ohne Richtung: c) Gesamtkraft 0 $\to$ 4 von 5\,P.}
\end{lpaufgabe}

\begin{lpaufgabe}{G4}{Aufzug fährt nach unten an: $72\;\text{kg}$, in $1.5\;\text{s}$ auf $2.7\;\text{m/s}$}{5 P}
\begin{raster}
Beschleunigung & 1\E & $a = \dfrac{\Delta v}{\Delta t} = \dfrac{2.7\;\text{m/s}}{1.5\;\text{s}} = 1.8\;\text{m/s}^2$ nach unten. Gleichwertig: $a = -1.8\;\text{m/s}^2$ mit «nach oben positiv». Ein Betrag ohne Richtung und ohne Vorzeichen gibt diesen Punkt nicht.\\
Ansatz & 1\W & nach oben positiv: $F_N - F_G = m \cdot a$, also $F_N = m \cdot (g + a)$ mit $a < 0$ (gleichwertig: $F_N = m \cdot (g - 1.8\;\text{m/s}^2)$)\\
Normalkraft & 1\E & $F_N = 72\;\text{kg} \cdot (9.81 - 1.8)\;\text{m/s}^2 \approx 577\;\text{N}$\\
Anzeige & 1\E & $\dfrac{F_N}{g} = \dfrac{577\;\text{N}}{9.81\;\text{m/s}^2} \approx 58.8\;\text{kg}$\\
andere Situation & 1\B & Die Fahrt nach oben wird mit $1.8\;\text{m/s}^2$ gebremst: Auch dann zeigt die Beschleunigung nach unten (Begründung nötig)\\
\end{raster}
\fehler{Ein Fehler, ein Abzug — jeder der beiden Ansatzfehler kostet einen Punkt, die Folgezeilen zählen:
Vorzeichen falsch, $F_N = 72\;\text{kg} \cdot (9.81 + 1.8)\;\text{m/s}^2 \approx 836\;\text{N}$, Anzeige $85.2\;\text{kg}$: Ansatz 0; Normalkraft und Anzeige folgerichtig $\to$ 4 von 5\,P.
Gewichtskraft vergessen, $F_N = m \cdot a \approx 130\;\text{N}$, Anzeige $13.2\;\text{kg}$: Ansatz 0; Normalkraft und Anzeige folgerichtig $\to$ 4 von 5\,P.
«Fährt nach unten mit konstantem Tempo»: andere Situation 0 $\to$ 4 von 5\,P.}
\end{lpaufgabe}

\begin{lpaufgabe}{G5}{Hängender Körper $m_2 = 0.5\;\text{kg}$, gemessen $F_S = 3.6\;\text{N}$}{4 P}
\begin{raster}
Ansatz & 1\W & am hängenden Körper (nach unten positiv): $m_2 \cdot g - F_S = m_2 \cdot a$\\
Beschleunigung & 1\E & $a = \dfrac{m_2 \cdot g - F_S}{m_2} = \dfrac{0.5\;\text{kg} \cdot 9.81\;\text{m/s}^2 - 3.6\;\text{N}}{0.5\;\text{kg}} \approx 2.61\;\text{m/s}^2$\\
Wagenmasse & 1\E & am Wagen allein: $F_S = m_1 \cdot a$, also $m_1 = \dfrac{3.6\;\text{N}}{2.61\;\text{m/s}^2} \approx 1.38\;\text{kg}$\\
doppelter Wagen & 1\B & Nein, mehr als halb so gross: Der Antrieb $m_2 \cdot g$ bleibt, die Gesamtmasse wächst nur von rund $1.88\;\text{kg}$ auf $3.26\;\text{kg}$, nicht auf das Doppelte ($a \approx 1.51\;\text{m/s}^2$ statt $1.31\;\text{m/s}^2$; die Rechnung ist nicht nötig)\\
\end{raster}
\fehler{$F_S = m_2 \cdot a$, also $a = 7.2\;\text{m/s}^2$ (Gewichtskraft am hängenden Körper vergessen): Ansatz 0; Beschleunigung 0; $m_1 = \dfrac{3.6\;\text{N}}{7.2\;\text{m/s}^2} = 0.5\;\text{kg}$ folgerichtig $\to$ 2 von 4\,P.
$m_1 = \dfrac{F_S}{g} \approx 0.37\;\text{kg}$ (Fadenkraft als Gewichtskraft des Wagens): Wagenmasse 0 $\to$ 3 von 4\,P.
«Ja, doppelte Masse, halbe Beschleunigung»: doppelter Wagen 0 $\to$ 3 von 4\,P.}
\end{lpaufgabe}

\begin{lpaufgabe}{G6}{Auto in der Kurve: $1200\;\text{kg}$, $72\;\text{km/h}$, Haftreibung $4000\;\text{N}$}{4 P}
\begin{raster}
Ansatz & 1\W & $F_z = \dfrac{m \cdot v^2}{r}$, also $r = \dfrac{m \cdot v^2}{F_z}$; dazu $v = \dfrac{72}{3.6}\;\text{m/s} = 20\;\text{m/s}$\\
Radius & 1\E & $r = \dfrac{1200\;\text{kg} \cdot (20\;\text{m/s})^2}{4000\;\text{N}} = 120\;\text{m}$\\
Rolle, Richtung & 1\B & Die Haftreibung übernimmt die Rolle der Zentripetalkraft; sie zeigt zur Kurvenmitte (beides nötig)\\
Glatteis & 1\B & Die Kurve braucht $4000\;\text{N}$ zur Mitte, die Haftreibung schafft höchstens $1500\;\text{N}$: Das Auto kommt nicht mehr auf den Kreis, es rutscht aus der Kurve und bewegt sich geradeaus weiter (tangential). «Wird nach aussen geschleudert» zählt nur mit diesem Grund.\\
\end{raster}
\fehler{Nicht in m/s umgerechnet, $r = 1555\;\text{m}$: Ansatz 0; Radius folgerichtig zählt $\to$ 3 von 4\,P.
Quadrat vergessen, $r = 6\;\text{m}$: Ansatz 0; Radius folgerichtig zählt $\to$ 3 von 4\,P.
«Die Zentrifugalkraft wird zu gross»: Glatteis 0 $\to$ 3 von 4\,P.}
\end{lpaufgabe}

\needspace{14\baselineskip}
\section*{4 · Selbsteinschätzung}
\begin{tabularx}{\linewidth}{@{}l X@{}}\toprule
22 – 25 P & Die geprüften Teile sitzen. Wo du Punkte verloren hast: das Kapitel dieser Aufgabe nochmals (Zuordnung unten).\\
17 – 21 P & Den schwächsten Teil nochmals: Simulation und Übungen des Kapitels, in dem du die meisten Punkte verloren hast.\\
11 – 16 P & Zurück zu den Kapiteln aller Aufgaben, in denen du Punkte verloren hast.\\
0 – 10 P & Zurück zu Kapitel 1 und von dort der Reihe nach weiter.\\\midrule
\multicolumn{2}{@{}p{\linewidth}@{}}{\small\textbf{Aufgabe $\to$ Kapitel} (gilt für alle Bereiche): G1 $\to$ Kapitel 1, 2 (K1) · G2, G3 $\to$ Kapitel 1, 2 (K1, K2) · G4 $\to$ Kapitel 3 (K2) · G5 $\to$ Kapitel 4 (K2) · G6 $\to$ Kapitel 5 (K2)}\\\bottomrule
\end{tabularx}
\end{document}
